\section{积分论与收敛定理}

\label{sec:2-2}

测度只给出了“集合有多大”，真正进入优化与概率时，我们关心的是函数在这些集合上的平均、总量与极限行为。Lebesgue 积分之所以在本书中不可替代，并不只是因为它比 Riemann 积分更一般，而是因为它允许我们在极限过程里保持对可积性的控制。本节的核心，就是把“积分是什么”和“什么时候可以把极限推进积分号内”两件事组织成一条可反复调用的技术链。

\subsection{Lebesgue 积分的构造}

\begin{definition}[简单函数与非负积分]\label{def:simple-function-integral}
设 $(\Omega,\mathscr F,\mu)$ 为测度空间。若函数 $\varphi\colon \Omega\to [0,+\infty)$ 可写成
\[
\varphi=\sum_{k=1}^{m} a_k \mathbf 1_{A_k},
\]
其中 $a_k\ge 0$，$A_k\in \mathscr F$ 两两不交，则称 $\varphi$ 为一个\emph{非负简单函数}，并定义
\[
\int_\Omega \varphi\,d\mu:=\sum_{k=1}^{m} a_k\,\mu(A_k).
\]
\end{definition}

\begin{definition}[Lebesgue 积分]\label{def:lebesgue-integral}
设 $f\colon \Omega\to [0,+\infty]$ 为非负可测函数。定义
\[
\int_\Omega f\,d\mu
:=
\sup\Bigl\{\int_\Omega \varphi\,d\mu:\ 0\le \varphi\le f,\ \varphi\text{ 为简单函数}\Bigr\}.
\]
对一般实值可测函数 $f$，写
\[
f^+=\max\{f,0\},\qquad f^-=\max\{-f,0\}.
\]
若
\[
\int_\Omega f^+\,d\mu<\infty,\qquad \int_\Omega f^-\,d\mu<\infty,
\]
则称 $f$ \emph{可积}，并定义
\[
\int_\Omega f\,d\mu:=\int_\Omega f^+\,d\mu-\int_\Omega f^-\,d\mu.
\]
\end{definition}

\paragraph{为什么强调 $L^1$}

\begin{remark}
定义~\ref{def:lebesgue-integral} 给出的可积函数类正是 $L^1(\mu)$ 的底层对象。后文讨论期望目标、风险泛函、弱收敛下的积分连续性时，真正起作用的往往不是函数的逐点值，而是它在 $L^1$ 中的等价类与积分控制。
\end{remark}

\begin{proposition}[积分的基本性质]\label{prop:integral-basic-properties}
若 $f,g$ 可积，$a,b\in \mathbb R$，则：
\begin{enumerate}
  \item $af+bg$ 可积，且
  \[
  \int (af+bg)\,d\mu=a\int f\,d\mu+b\int g\,d\mu;
  \]
  \item 若 $f\le g$ a.e.，则
  \[
  \int f\,d\mu\le \int g\,d\mu;
  \]
  \item
  \[
  \left|\int f\,d\mu\right|\le \int |f|\,d\mu.
  \]
\end{enumerate}
\end{proposition}

\begin{proof}
对简单函数结论显然，由上确界定义与单调逼近推广到非负可测函数；再对一般可积函数分解为正负部分即可。第三条由
\[
-|f|\le f\le |f|
\]
和第二条直接推出。
\end{proof}

\subsection{单调收敛、Fatou 引理与控制收敛}

\begin{theorem}[单调收敛定理]\label{thm:monotone-convergence}
设 $\{f_n\}_{n\ge1}$ 为非负可测函数列，且
\[
0\le f_1\le f_2\le \cdots,\qquad f_n\uparrow f \text{ a.e.}
\]
则
\[
\int_\Omega f_n\,d\mu \uparrow \int_\Omega f\,d\mu.
\]
\end{theorem}

\begin{proof}
由单调性先得
\[
\int f_n\,d\mu\le \int f\,d\mu,
\]
故左侧极限存在且不超过右侧。设
\[
L:=\sup_n \int f_n\,d\mu.
\]
任取简单函数 $\varphi$ 满足 $0\le \varphi\le f$。对任意 $\varepsilon\in(0,1)$，考虑集合
\[
E_n:=\{\,\omega:\ f_n(\omega)\ge (1-\varepsilon)\varphi(\omega)\,\}.
\]
由于 $f_n\uparrow f\ge \varphi$，故 $\mathbf 1_{E_n}\uparrow 1$。于是由简单函数积分与测度从下连续性可得
\[
\int (1-\varepsilon)\varphi\,d\mu
=
\lim_{n\to\infty}\int_{E_n}(1-\varepsilon)\varphi\,d\mu
\le
\limsup_{n\to\infty}\int f_n\,d\mu
=L.
\]
令 $\varepsilon\downarrow 0$，再对所有 $\varphi\le f$ 取上确界，得到
\[
\int f\,d\mu\le L.
\]
结合前面的反向不等式，结论成立。
\end{proof}

\begin{lemma}[Fatou 引理]\label{lem:fatou}
设 $\{f_n\}_{n\ge1}$ 为非负可测函数列，则
\[
\int_\Omega \liminf_{n\to\infty} f_n\,d\mu
\le
\liminf_{n\to\infty}\int_\Omega f_n\,d\mu.
\]
\end{lemma}

\begin{proof}
设
\[
g_n:=\inf_{k\ge n} f_k.
\]
则 $\{g_n\}$ 仍为非负可测函数列，且 $g_n\uparrow \liminf_{n\to\infty}f_n$。由定理~\ref{thm:monotone-convergence}，
\[
\int \liminf_{n\to\infty}f_n\,d\mu
=
\lim_{n\to\infty}\int g_n\,d\mu.
\]
又对每个固定 $n$ 与任意 $k\ge n$，都有 $g_n\le f_k$，故
\[
\int g_n\,d\mu\le \inf_{k\ge n}\int f_k\,d\mu.
\]
令 $n\to\infty$ 即得结论。
\end{proof}

\begin{theorem}[控制收敛定理]\label{thm:dominated-convergence}
设 $\{f_n\}_{n\ge1}$ 为可测函数列，且
\[
f_n\to f \text{ a.e.},\qquad |f_n|\le g \text{ a.e.},
\]
其中 $g$ 可积。则 $f$ 可积，且
\[
\int_\Omega f_n\,d\mu\to \int_\Omega f\,d\mu.
\]
\end{theorem}

\begin{proof}
由 $|f_n|\le g$ 和逐点收敛可得 $|f|\le g$ a.e.，故 $f$ 可积。对非负函数列 $g+f_n$ 应用 Fatou 引理，得
\[
\int (g+f)\,d\mu
\le
\liminf_{n\to\infty}\int (g+f_n)\,d\mu
=
\int g\,d\mu+\liminf_{n\to\infty}\int f_n\,d\mu.
\]
从而
\[
\int f\,d\mu\le \liminf_{n\to\infty}\int f_n\,d\mu.
\]
同理，对非负函数列 $g-f_n$ 应用 Fatou 引理，得到
\[
\int (g-f)\,d\mu
\le
\liminf_{n\to\infty}\int (g-f_n)\,d\mu
=
\int g\,d\mu-\limsup_{n\to\infty}\int f_n\,d\mu,
\]
即
\[
\limsup_{n\to\infty}\int f_n\,d\mu\le \int f\,d\mu.
\]
结合二者便得所求极限。
\end{proof}

\begin{proposition}[积分的绝对连续性]\label{prop:absolute-continuity-integral}
若 $f\in L^1(\mu)$，则对任意 $\varepsilon>0$，存在 $\delta>0$，使得只要 $A\in \mathscr F$ 且 $\mu(A)<\delta$，就有
\[
\int_A |f|\,d\mu<\varepsilon.
\]
\end{proposition}

\begin{proof}
取 $M>0$ 使
\[
\int_{\{|f|>M\}} |f|\,d\mu<\frac{\varepsilon}{2}.
\]
再令 $\delta=\varepsilon/(2M)$。若 $\mu(A)<\delta$，则
\[
\int_A |f|\,d\mu
=
\int_{A\cap\{|f|\le M\}}|f|\,d\mu+\int_{A\cap\{|f|>M\}}|f|\,d\mu
\le
M\mu(A)+\frac{\varepsilon}{2}
<
\varepsilon.
\]
\end{proof}

\begin{example}[抽象例子：指标函数逼近]\label{ex:indicator-approximation}
设 $A_n\uparrow A$，则 $\mathbf 1_{A_n}\uparrow \mathbf 1_A$。由定理~\ref{thm:monotone-convergence} 得
\[
\mu(A_n)=\int \mathbf 1_{A_n}\,d\mu \uparrow \int \mathbf 1_A\,d\mu=\mu(A).
\]
把这个例子放在这里，是为了说明测度从下连续性其实是积分理论中最基本的单调收敛现象。
\end{example}

\begin{example}[有限维坍缩：随机向量的期望]\label{ex:finite-dimensional-expectation}
若 $X$ 是 $\mathbb R^n$ 值随机向量，$h\colon \mathbb R^n\to \mathbb R$ 可积，则
\[
\mathbb E[h(X)]=\int_{\mathbb R^n} h(x)\,\mu_X(dx).
\]
把这个例子放在这里，是为了提醒读者 Boyd 书里常见的“期望目标函数”并不是一种独立算子，而只是定义~\ref{def:lebesgue-integral} 在概率测度上的写法。
\end{example}

\begin{example}[应用触点：Monte Carlo 与风险泛函]\label{ex:monte-carlo-risk}
设损失函数 $L_n$ 由模拟近似得到，且 $L_n\to L$ a.s.，同时存在可积上界 $G$ 使 $|L_n|\le G$。由定理~\ref{thm:dominated-convergence}，
\[
\mathbb E[L_n]\to \mathbb E[L].
\]
把这个例子放在这里，是为了说明后文量化风险优化和分布鲁棒目标中，控制收敛定理是“极限可推进期望号内”的最基本合法性来源。
\end{example}

\subsection*{有限维对照}

有限维教材里，人们经常把“对随机变量取期望”“把极限与期望交换”写成直观计算；而在无穷维和函数空间中，这些动作是否成立完全依赖定理~\ref{thm:monotone-convergence}、引理~\ref{lem:fatou} 和定理~\ref{thm:dominated-convergence} 所提供的条件链。

\subsection*{习题（待补）}

\begin{exercise}[单调收敛占位]\label{exr:monotone-convergence-placeholder}
补一题：用定理~\ref{thm:monotone-convergence} 证明概率测度对递增事件列的从下连续性。
\end{exercise}

\begin{exercise}[控制收敛桥接题占位]\label{exr:dominated-convergence-placeholder}
补一题：把例~\ref{ex:monte-carlo-risk} 落回欧氏空间中的随机向量积分，说明控制上界在数值优化中为何常以矩条件或截断条件的形式出现。
\end{exercise}
